\begin{textsample}{POS Dim 1 – 2013 – Score 19.00 – 2013/t001983.txt}  \label{ex:f1_pos_044}
When I was a young man, I spent six years of wild adventure in \textbf{the} tropics working as an investigative journalist in some of \textbf{the} most bewitching parts of \textbf{the} world.
I was as reckless and foolish as only young men can be.
This is why wars get fought.
But I also felt more alive than I ' ve ever done since.
And when I came home, I found \textbf{the} scope of my existence gradually diminishing until loading \textbf{the} dishwasher seemed like an interesting challenge.
And I found myself sort of scratching at \textbf{the} walls of life, as if I was trying to find a way out into a wider space beyond.
I was, I believe, ecologically bored.
Now, we \textbf{evolved} in rather more challenging times than these, in a world of horns and tusks and fangs and claws.
And we still possess \textbf{the} fear and \textbf{the} courage and \textbf{the} aggression required to navigate those times.
But in our comfortable, safe, crowded lands, we have few opportunities to exercise them without harming other people.
And this was \textbf{the} sort of constraint that I found myself bumping up against.
To conquer uncertainty, to know what comes next, that ' s almost been \textbf{the} dominant aim of industrialized societies, and having got there, or almost got there, we have just encountered a new set of unmet needs.
We ' ve privileged safety over experience and we ' ve gained a lot in doing so, but I think we ' ve lost something too.
Now, I don ' t romanticize \textbf{evolutionary} time.
I ' m already beyond \textbf{the} lifespan of most hunter-gatherers, and \textbf{the} outcome of a mortal combat between me myopically stumbling around with a stone-tipped spear and an enraged giant aurochs isn ' t very hard to predict.
Nor was it authenticity that I was looking for.
I don ' t find that a useful or even intelligible concept.
I just wanted a richer and rawer life than I ' ve been able to lead in Britain, or, indeed, that we can lead in most parts of \textbf{the} industrialized world.
And it was only when I stumbled across an unfamiliar word that I began to understand what I was looking for.
And as soon as I found that word, I realized that I wanted to devote much of \textbf{the} rest of my life to it. \textbf{The} word is"rewilding,"and even though rewilding is a young word, it already has several definitions.
But there are two in particular that fascinate me. \textbf{The} first one is \textbf{the} mass restoration of ecosystems.
One of \textbf{the} most exciting scientific findings of \textbf{the} past half century has been \textbf{the} discovery of widespread trophic cascades.
A trophic cascade is an ecological process which starts at \textbf{the} top of \textbf{the} food chain and tumbles all \textbf{the} way down to \textbf{the} bottom, and \textbf{the} classic example is what happened in \textbf{the} Yellowstone National Park in \textbf{the} United States when wolves were reintroduced in 1995.
Now, we all know that wolves kill various species of animals, but perhaps we ' re slightly less aware that they give life to many others.
It sounds strange, but just follow me for a while.
Before \textbf{the} wolves turned up, they ' d been absent for 70 years. \textbf{The} numbers of deer, because there was nothing to hunt them, had built up and built up in \textbf{the} Yellowstone Park, and despite efforts by humans to control them, they ' d managed to reduce much of \textbf{the} vegetation there to almost nothing, they ' d just grazed it away.
But as soon as \textbf{the} wolves arrived, even though they were few in number, they started to have \textbf{the} most remarkable effects.
First, of course, they killed some of \textbf{the} deer, but that wasn ' t \textbf{the} major thing.
Much more significantly, they radically changed \textbf{the} behavior of \textbf{the} deer. \textbf{The} deer started avoiding certain parts of \textbf{the} park, \textbf{the} places where they could be trapped most easily, particularly \textbf{the} valleys and \textbf{the} gorges, and immediately those places started to regenerate.
In some areas, \textbf{the} height of \textbf{the} trees quintupled in just six years.
Bare valley sides quickly became forests of aspen and willow and cottonwood.
And as soon as that happened, \textbf{the} birds started moving in. \textbf{The} number of songbirds, of migratory birds, started to increase greatly. \textbf{The} number of beavers started to increase, because beavers like to eat \textbf{the} trees.
And beavers, like wolves, are ecosystem engineers.
They create niches for other species.
And \textbf{the} dams they built in \textbf{the} rivers provided \textbf{habitats} for otters and muskrats and ducks and \textbf{fish} and reptiles and amphibians. \textbf{The} wolves killed coyotes, and as a result of that, \textbf{the} number of rabbits and mice began to rise, which meant more hawks, more weasels, more foxes, more badgers.
Ravens and bald eagles came down to feed on \textbf{the} carrion that \textbf{the} wolves had left.
Bears fed on it too, and their population began to rise as well, partly also because there were more berries growing on \textbf{the} regenerating shrubs, and \textbf{the} bears reinforced \textbf{the} impact of \textbf{the} wolves by killing some of \textbf{the} calves of \textbf{the} deer.
But here ' s where it gets really interesting. \textbf{The} wolves changed \textbf{the} behavior of \textbf{the} rivers.
They began to meander less.
There was less erosion. \textbf{The} channels narrowed.
More pools formed, more riffle sections, all of which were great for wildlife \textbf{habitats}. \textbf{The} rivers changed in response to \textbf{the} wolves, and \textbf{the} reason was that \textbf{the} regenerating forests stabilized \textbf{the} banks so that they collapsed less often, so that \textbf{the} rivers became more fixed in their course.
Similarly, by driving \textbf{the} deer out of some places and \textbf{the} vegetation recovering on \textbf{the} valley sides, there was less soil erosion, because \textbf{the} vegetation stabilized that as well.
So \textbf{the} wolves, small in number, transformed not just \textbf{the} ecosystem of \textbf{the} Yellowstone National Park, this huge area of land, but also its physical geography. \textbf{Whales} in \textbf{the} southern oceans have similarly wide-ranging effects.
One of \textbf{the} many post-rational excuses made by \textbf{the} Japanese government for killing \textbf{whales} is that they said,"Well, \textbf{the} number of \textbf{fish} and krill will rise and then there ' ll be more for people to eat."Well, it ' s a stupid excuse, but it sort of kind of makes sense, doesn ' t it, because you ' d think that \textbf{whales} eat huge amounts of \textbf{fish} and krill, so obviously take \textbf{the} \textbf{whales} away, there ' ll be more \textbf{fish} and krill.
But \textbf{the} opposite happened.
You take \textbf{the} \textbf{whales} away, and \textbf{the} number of krill collapses.
Why would that possibly have happened?
Well, it now turns out that \textbf{the} \textbf{whales} are crucial to sustaining that entire ecosystem, and one of \textbf{the} reasons for this is that they often feed at depth and then they come up to \textbf{the} \textbf{surface} and produce what biologists politely call large fecal plumes, huge \textbf{explosions} of poop right across \textbf{the} \textbf{surface} waters, up in \textbf{the} photic zone, where there ' s enough light to allow photosynthesis to take place, and those great plumes of fertilizer stimulate \textbf{the} growth of phytoplankton, \textbf{the} plant plankton at \textbf{the} bottom of \textbf{the} food chain, which stimulate \textbf{the} growth of zooplankton, which feed \textbf{the} \textbf{fish} and \textbf{the} krill and all \textbf{the} rest of it. \textbf{The} other thing that \textbf{whales} do is that, as they ' re plunging up and down through \textbf{the} water column, they ' re kicking \textbf{the} phytoplankton back up towards \textbf{the} \textbf{surface} where it can continue to survive and reproduce.
And interestingly, well, we know that plant plankton in \textbf{the} oceans absorb \textbf{carbon} from \textbf{the} atmosphere — \textbf{the} more plant plankton there are, \textbf{the} more \textbf{carbon} they absorb — and eventually they filter down into \textbf{the} \textbf{abyss} and remove that \textbf{carbon} from \textbf{the} atmospheric system.
Well, it seems that when \textbf{whales} were at their historic populations, they were probably responsible for sequestering some tens of millions of tons of \textbf{carbon} every year from \textbf{the} atmosphere.
And when you look at it like that, you think, wait a minute, here are \textbf{the} wolves changing \textbf{the} physical geography of \textbf{the} Yellowstone National Park.
Here are \textbf{the} \textbf{whales} changing \textbf{the} composition of \textbf{the} atmosphere.
You begin to see that possibly, \textbf{the} evidence supporting James Lovelock ' s Gaia hypothesis, which conceives of \textbf{the} world as a coherent, self-regulating organism, is beginning, at \textbf{the} ecosystem level, to accumulate.
Trophic cascades tell us that \textbf{the} natural world is even more fascinating and complex than we thought it was.
They tell us that when you take away \textbf{the} large animals, you are left with a radically different ecosystem to one which retains its large animals.
And they make, in my view, a powerful case for \textbf{the} reintroduction of missing species.
Rewilding, to me, means bringing back some of \textbf{the} missing plants and animals.
It means taking down \textbf{the} fences, it means blocking \textbf{the} drainage ditches, it means preventing commercial \textbf{fishing} in some large areas of \textbf{sea}, but otherwise stepping back.
It has no view as to what a right ecosystem or a right assemblage of species looks like.
It doesn ' t try to produce a heath or a meadow or a rain forest or a kelp garden or a coral reef.
It lets nature decide, and nature, by and large, is pretty good at deciding.
Now, I mentioned that there are two definitions of rewilding that interest me. \textbf{The} other one is \textbf{the} rewilding of human life.
And I don ' t see this as an alternative to civilization.
I believe we can \textbf{enjoy} \textbf{the} benefits of advanced technology, as we ' re doing now, but at \textbf{the} same time, if we choose, have access to a richer and wilder life of adventure when we want to because there would be wonderful, rewilded \textbf{habitats}.
And \textbf{the} opportunities for this are developing more rapidly than you might think possible.
There ' s one estimate which suggests that in \textbf{the} United States, two thirds of \textbf{the} land which was once forested and then cleared has become reforested as loggers and farmers have retreated, particularly from \textbf{the} \textbf{eastern} half of \textbf{the} country.
There ' s another one which suggests that 30 million hectares of land in Europe, an area \textbf{the} size of Poland, will be vacated by farmers between 2000 and 2030.
Now, faced with opportunities like that, does it not seem a little unambitious to be thinking only of bringing back wolves, lynx, bears, beavers, bison, boar, moose, and all \textbf{the} other species which are already beginning to move quite rapidly across Europe?
Perhaps we should also start thinking about \textbf{the} return of some of our lost megafauna.
What megafauna, you say?
Well, every \textbf{continent} had one, apart from Antarctica.
When Trafalgar Square in London was excavated, \textbf{the} river gravels there were found to be stuffed with \textbf{the} bones of hippopotamus, rhinos, \textbf{elephants}, hyenas, \textbf{lions}.
Yes, ladies and gentlemen, there were \textbf{lions} in Trafalgar Square long before Nelson ' s Column was built.
All these species lived here in \textbf{the} last interglacial period, when temperatures were pretty similar to our own.
It ' s not climate, largely, which has got rid of \textbf{the} world ' s megafaunas.
It ' s pressure from \textbf{the} human population hunting and destroying their \textbf{habitats} which has done so.
And even so, you can still see \textbf{the} shadows of these great beasts in our current ecosystems.
Why is it that so many deciduous trees are able to sprout from whatever point \textbf{the} trunk is broken?
Why is it that they can withstand \textbf{the} loss of so much of their bark?
Why do understory trees, which are subject to lower \textbf{sheer} forces from \textbf{the} wind and have to carry less weight than \textbf{the} big canopy trees, why are they so much tougher and harder to break than \textbf{the} canopy trees are? \textbf{Elephants}.
They are elephant-adapted.
In Europe, for example, they \textbf{evolved} to resist \textbf{the} straight-tusked \textbf{elephant}, elephas antiquus, which was a great beast.
It was related to \textbf{the} Asian \textbf{elephant}, but it was a temperate animal, a temperate forest creature.
It was a lot bigger than \textbf{the} Asian \textbf{elephant}.
But why is it that some of our common shrubs have spines which seem to be over- engineered to resist browsing by deer?
Perhaps because they \textbf{evolved} to resist browsing by rhinoceros.
Isn ' t it an amazing thought that every time you wander into a park or down an avenue or through a leafy street, you can see \textbf{the} shadows of these great beasts?
Paleoecology, \textbf{the} study of past ecosystems, crucial to an understanding of our own, feels like a portal through which you may pass into an enchanted kingdom.
And if we really are looking at areas of land of \textbf{the} sort of sizes I ' ve been talking about becoming available, why not reintroduce some of our lost megafauna, or at least species closely related to those which have become extinct everywhere?
Why shouldn ' t all of us have a Serengeti on our doorsteps?
And perhaps this is \textbf{the} most important thing that rewilding offers us, \textbf{the} most important thing that ' s missing from our lives : hope.
In motivating people to love and defend \textbf{the} natural world, an ounce of hope is worth a ton of despair. \textbf{The} story rewilding tells us is that ecological change need not always proceed in one direction.
It offers us \textbf{the} hope that our silent \textbf{spring} could be replaced by a raucous summer.
Thank you.

% matched lemmas: abyss, carbon, continent, eastern, elephant, enjoy, evolutionary, evolve, explosion, fish, fishing, habitat, lion, sea, sheer, spring, surface, the, whale
\end{textsample}
